\documentclass[11pt]{article}

% -------------------- Packages --------------------
\usepackage[a4paper,margin=1in]{geometry}
\usepackage{hyperref}
\usepackage{enumitem}
\usepackage{titlesec}
\usepackage{fancyhdr}
\usepackage{xcolor}
\usepackage{listings}
\usepackage{courier}

% -------------------- Hyperref Setup --------------------
\hypersetup{
  colorlinks=true,
  linkcolor=blue,
  urlcolor=blue,
  citecolor=blue
}

% -------------------- Header/Footer --------------------
\pagestyle{fancy}
\fancyhf{}
\lhead{MCV lab}
\rhead{Scalable ML Experimentation}
\cfoot{\thepage}

% -------------------- Section Formatting --------------------
\titleformat{\section}{\Large\bfseries}{\thesection}{0.8em}{}
\titleformat{\subsection}{\normalsize\bfseries}{\thesubsection}{0.8em}{}

\titlespacing*{\section}{0pt}{1.5em}{0.8em}
\titlespacing*{\subsection}{0pt}{1.0em}{0.5em}

% -------------------- Global List Formatting --------------------
\setlist[itemize]{
  label={-},
  leftmargin=2em,
  labelsep=0.6em,
  itemsep=0.2em
}
\setlist[itemize,2]{itemsep=0.15em}
\setlist[enumerate]{leftmargin=2em,itemsep=0.15em}

% -------------------- Listings / Code Style --------------------
\definecolor{codegray}{rgb}{0.95,0.95,0.95}
\definecolor{commentgreen}{rgb}{0.0,0.5,0.0}

\lstset{
  basicstyle=\footnotesize\ttfamily,
  backgroundcolor=\color{codegray},
  frame=single,
  breaklines=true,
  showstringspaces=false,
  keywordstyle=\color{blue},
  commentstyle=\color{commentgreen},
  stringstyle=\color{red},
  columns=fullflexible,
  keepspaces=true,
  tabsize=4,
  xleftmargin=0.02\textwidth,
  xrightmargin=0.02\textwidth
}

% % -------------------- Document Metadata --------------------
% \title{\textbf{Lab Standard: Scalable ML Experimentation with Argument Parsers}}
% \author{Medical Computer Vision Lab (MCV Lab)}
% \date{\today}

% \begin{document}
% \maketitle

% \tableofcontents
% \newpage

% -------------------- Document Metadata --------------------
\title{\textbf{Medical Computer Vision Lab \\ \small{
Scalable ML Experimentation with Argument Parsers: Lab Standard}}}
\author{Maintainer: Ashery Mbilinyi}
\date{Version 1.0 : \today}

\begin{document}
\maketitle

\tableofcontents
\newpage



% % -------------------- Abstract (new page + included in ToC) --------------------
% \addcontentsline{toc}{section}{Abstract}
% \begin{abstract}
% This document defines our lab standard for writing scalable machine learning (ML) experiments code using argument parsers.
% It ensures experiments are reproducible, auditable, scalable, and easy to run by new lab members without modifying source code.
% \end{abstract}
% \newpage

\section{Purpose}
This document defines our lab standard for writing scalable machine learning (ML) experiment code using argument parsers.
The goal is to ensure our experiments are:
\begin{itemize}
    \item \textbf{Reproducible}: every run can be reproduced exactly
    \item \textbf{Auditable}: every result is traceable to a configuration
    \item \textbf{Scalable}: easy to run hundreds of experiments on local machines, servers, and HPC clusters
    \item \textbf{Collaborative}: new students can run baselines without modifying code
\end{itemize}

\section{Why Argument Parsers Enable Scalability}
Argument parsers (e.g., Python \texttt{argparse}, Hydra, Typer) improve scalability by turning scripts into \textbf{configurable tools} rather than one-off code.

\subsection{Key Benefits}
\begin{itemize}
    \item \textbf{\small Eliminates code editing for experimentation}.\\
    Parameters should not require editing the training script. Instead, experiments are configured from the command line.

    \item \textbf{Reproducibility at scale}.\\
    Each run can save its full configuration (\texttt{config.json}, \texttt{cmd.txt}), enabling exact reproduction months later.

    \item \textbf{Simple sweeps and automation}.\\
    Sweeps become trivial through bash loops, schedulers, or frameworks (SLURM, Ray Tune, W\&B sweeps).

    \item \textbf{HPC and pipeline compatibility}.\\
    Cluster job submission is easiest when scripts take structured arguments and output to predictable directories.

    \item \textbf{Safer experiments}.\\
    Parsers support type-checking, required fields, and limited choices, reducing wasted GPU runs.
\end{itemize}

\section{Lab Policy: No Code Changes for Experiments}

\subsection{Rule}
\textbf{Model experimentation must not require changing source code.}

\subsection{Allowed Ways to Configure Experiments}
\begin{itemize}
    \item CLI arguments (\texttt{--lr}, \texttt{--batch\_size}, \texttt{--seed}, etc.)
    \item Config files (\texttt{--config configs/baseline.json})
\end{itemize}

\subsection{Disallowed}
\begin{itemize}
    \item Editing constants inside scripts (learning rate, batch size, epochs)
    \item Hardcoding dataset paths in code
    \item Changing the script for each experiment
\end{itemize}

\section{Standard Experiment Output Structure}
Every experiment run should create a unique run directory containing:
\begin{itemize}
    \item \texttt{config.json} : resolved configuration used for the run
    \item \texttt{cmd.txt} : command used to launch the run
    \item \texttt{checkpoints/} : saved model checkpoints
    \item \texttt{logs/} or metrics files
\end{itemize}

Example:
\begin{lstlisting}
runs/
  rsna_trauma_resnet50_lr0.0001_20260110_215000/
    config.json
    cmd.txt
    checkpoints/
      last.ckpt
      best.ckpt
\end{lstlisting}

\section{Core Arguments Required in All Training Scripts}
At minimum, every training script must support:

\subsection{Experiment / Logging}
\begin{itemize}
    \item \texttt{--run\_name} (auto-generated if empty)
    \item \texttt{--out\_dir} (default: \texttt{runs})
\end{itemize}

\subsection{Dataset}
\begin{itemize}
    \item \texttt{--dataset} (choices must be enforced)
    \item \texttt{--data\_dir}
    \item \texttt{--num\_workers}
\end{itemize}

\subsection{Model}
\begin{itemize}
    \item \texttt{--model} (choices enforced)
    \item \texttt{--pretrained}
\end{itemize}

\subsection{Optimization}
\begin{itemize}
    \item \texttt{--epochs}
    \item \texttt{--batch\_size}
    \item \texttt{--lr}
    \item \texttt{--weight\_decay}
    \item \texttt{--optimizer} (choices enforced)
    \item \texttt{--scheduler} (choices enforced)
\end{itemize}

\subsection{Reproducibility}
\begin{itemize}
    \item \texttt{--seed}
    \item \texttt{--deterministic}
\end{itemize}

\subsection{Hardware}
\begin{itemize}
    \item \texttt{--device} (cpu/cuda/mps)
    \item \texttt{--amp}
\end{itemize}

\subsection{Resume / Checkpointing}
\begin{itemize}
    \item \texttt{--resume} (checkpoint path)
\end{itemize}

\subsection{Config File Support (Mandatory)}
\begin{itemize}
    \item \texttt{--config} (JSON/YAML config)
\end{itemize}

\section{Standard CLI Usage Examples}

\subsection{Baseline run}
\begin{lstlisting}
python train.py --dataset rsna_trauma --model resnet50 --lr 1e-4 --seed 0 --amp
\end{lstlisting}

\subsection{Config-based run}
\begin{lstlisting}
python train.py --config configs/baseline.json
\end{lstlisting}

\subsection{Config run with overrides}
\begin{lstlisting}
python train.py --config configs/baseline.json --lr 5e-5 --seed 2
\end{lstlisting}

\subsection{Sweeps}
\begin{lstlisting}
for seed in 0 1 2; do
  for lr in 1e-3 1e-4 1e-5; do
    python train.py \
      --dataset rsna_trauma \
      --model resnet50 \
      --lr $lr \
      --seed $seed \
      --amp
  done
done
\end{lstlisting}

\section{Recommended Implementation Template}
Below is the lab standard \texttt{argparse} training entrypoint template. Students may adapt it, but required fields and output structure must remain consistent.

\subsection{Full Reference Template}
\begin{lstlisting}[language=Python]
#!/usr/bin/env python3
"""
train.py - scalable training entrypoint template

Usage:
  python train.py --dataset mimic_ecg --model resnet18 --lr 1e-4 --seed 0
  python train.py --config configs/baseline.json --batch_size 64
  python train.py --resume runs/exp_001/checkpoints/last.ckpt
"""

from __future__ import annotations

import argparse
import json
import random
import shutil
import sys
from dataclasses import asdict, dataclass
from datetime import datetime
from pathlib import Path
from typing import Any

import numpy as np

try:
    import torch
except ImportError:
    torch = None


def set_seed(seed: int, deterministic: bool = False) -> None:
    random.seed(seed)
    np.random.seed(seed)

    if torch is not None:
        torch.manual_seed(seed)
        torch.cuda.manual_seed_all(seed)

        if deterministic:
            torch.backends.cudnn.deterministic = True
            torch.backends.cudnn.benchmark = False


@dataclass
class TrainConfig:
    # Experiment / logging
    run_name: str
    out_dir: str
    save_every: int
    log_every: int

    # Dataset
    dataset: str
    data_dir: str
    num_workers: int

    # Model
    model: str
    pretrained: bool

    # Optimization
    epochs: int
    batch_size: int
    lr: float
    weight_decay: float
    optimizer: str
    scheduler: str

    # Reproducibility
    seed: int
    deterministic: bool

    # Hardware
    device: str
    amp: bool

    # Resume
    resume: str | None


def timestamp() -> str:
    return datetime.now().strftime("%Y%m%d_%H%M%S")


def ensure_dir(path: str | Path) -> Path:
    path = Path(path)
    path.mkdir(parents=True, exist_ok=True)
    return path


def dump_json(path: Path, obj: dict[str, Any]) -> None:
    with open(path, "w", encoding="utf-8") as f:
        json.dump(obj, f, indent=2, sort_keys=True)


def load_json(path: str | Path) -> dict[str, Any]:
    with open(path, "r", encoding="utf-8") as f:
        return json.load(f)


def save_run_config(run_dir: Path, cfg: TrainConfig, argv: list[str]) -> None:
    dump_json(run_dir / "config.json", asdict(cfg))
    (run_dir / "cmd.txt").write_text("python " + " ".join(argv) + "\n", encoding="utf-8")


def build_parser() -> argparse.ArgumentParser:
    parser = argparse.ArgumentParser(
        prog="train.py",
        description="Scalable training entrypoint",
        formatter_class=argparse.ArgumentDefaultsHelpFormatter,
    )

    g = parser.add_argument_group("Experiment / Logging")
    g.add_argument("--run_name", type=str, default="", help="Run name (auto if empty).")
    g.add_argument("--out_dir", type=str, default="runs", help="Base output directory.")
    g.add_argument("--save_every", type=int, default=1, help="Checkpoint every N epochs.")
    g.add_argument("--log_every", type=int, default=50, help="Log every N steps.")

    g = parser.add_argument_group("Dataset")
    g.add_argument(
        "--dataset",
        type=str,
        default="dummy",
        choices=["dummy", "ptbxl", "mimic_ecg", "rsna_trauma"],
    )
    g.add_argument("--data_dir", type=str, default="data")
    g.add_argument("--num_workers", type=int, default=4)

    g = parser.add_argument_group("Model")
    g.add_argument(
        "--model",
        type=str,
        default="resnet18",
        choices=["resnet18", "resnet50", "unet", "vit_b16"],
    )
    g.add_argument("--pretrained", action="store_true")

    g = parser.add_argument_group("Optimization")
    g.add_argument("--epochs", type=int, default=50)
    g.add_argument("--batch_size", type=int, default=32)
    g.add_argument("--lr", type=float, default=1e-4)
    g.add_argument("--weight_decay", type=float, default=1e-4)
    g.add_argument("--optimizer", type=str, default="adamw", choices=["sgd", "adam", "adamw"])
    g.add_argument("--scheduler", type=str, default="cosine", choices=["none", "step", "cosine"])

    g = parser.add_argument_group("Reproducibility")
    g.add_argument("--seed", type=int, default=0)
    g.add_argument("--deterministic", action="store_true",
                   help="Force deterministic behavior (slower).")

    g = parser.add_argument_group("Hardware")
    g.add_argument("--device", type=str, default="cuda", choices=["cuda", "cpu", "mps"])
    g.add_argument("--amp", action="store_true", help="Enable mixed precision.")

    g = parser.add_argument_group("Resume")
    g.add_argument("--resume", type=str, default=None)

    parser.add_argument("--config", type=str, default=None,
                        help="Optional JSON config to load defaults from.")

    return parser


def parse_args_with_config(parser: argparse.ArgumentParser, argv: list[str]) -> argparse.Namespace:
    partial, _ = parser.parse_known_args(argv)

    if partial.config is not None:
        cfg = load_json(partial.config)
        parser.set_defaults(**cfg)

    args = parser.parse_args(argv)
    return args


def train(cfg: TrainConfig, run_dir: Path) -> None:
    print(json.dumps(asdict(cfg), indent=2))
    set_seed(cfg.seed, deterministic=cfg.deterministic)

    # TODO: implement actual training loop
    print(f"Training would run here. Outputs at: {run_dir}")


def main(argv: list[str]) -> int:
    parser = build_parser()
    args = parse_args_with_config(parser, argv)

    if args.run_name.strip() == "":
        args.run_name = f"{args.dataset}_{args.model}_lr{args.lr}_{timestamp()}"

    base_out = Path(args.out_dir)
    run_dir = ensure_dir(base_out / args.run_name)
    ensure_dir(run_dir / "checkpoints")

    if args.config is not None:
        shutil.copy(args.config, run_dir / "config_used.json")

    cfg = TrainConfig(
        run_name=args.run_name,
        out_dir=str(run_dir),
        save_every=args.save_every,
        log_every=args.log_every,
        dataset=args.dataset,
        data_dir=args.data_dir,
        num_workers=args.num_workers,
        model=args.model,
        pretrained=args.pretrained,
        epochs=args.epochs,
        batch_size=args.batch_size,
        lr=args.lr,
        weight_decay=args.weight_decay,
        optimizer=args.optimizer,
        scheduler=args.scheduler,
        seed=args.seed,
        deterministic=args.deterministic,
        device=args.device,
        amp=args.amp,
        resume=args.resume,
    )

    save_run_config(run_dir, cfg, argv)
    train(cfg, run_dir)
    return 0


if __name__ == "__main__":
    raise SystemExit(main(sys.argv[1:]))
\end{lstlisting}

\section{Config File Example}
Example file: \texttt{configs/baseline.json}
\begin{lstlisting}
{
  "dataset": "rsna_trauma",
  "data_dir": "/datasets/rsna",
  "model": "resnet50",
  "epochs": 30,
  "batch_size": 16,
  "lr": 0.0001,
  "optimizer": "adamw",
  "scheduler": "cosine",
  "seed": 0,
  "amp": true
}
\end{lstlisting}

\section{Checklist for Code Reviews}
Before merging new training pipelines into the lab repo, ensure:
\begin{itemize}
    \item Training can run fully from CLI with no code edits
    \item \texttt{--config} works and CLI overrides it
    \item Every run saves \texttt{config.json} and \texttt{cmd.txt}
    \item Output directory structure matches lab standard
    \item Seed is set and reproducibility toggles exist
    \item Dataset/model/optimizer use \texttt{choices} where applicable
\end{itemize}

\section{Recommended Extensions (Optional)}
\begin{itemize}
    \item Add W\&B or MLflow logging of the full config
    \item Add validation that checks incompatible arguments
    \item Add support for YAML configs (OmegaConf)
    \item Upgrade to Hydra for hierarchical config management
\end{itemize}

\end{document}
